\documentclass{beamer}
\usepackage[utf8]{inputenc}

\usetheme{Berkeley}
\usecolortheme{default}

\setbeamertemplate{footline}[frame number]
\beamertemplatenavigationsymbolsempty


\title{Developing Effective Writing Habits}
\subtitle{Systems, Structure, and Sustainability}
\author[]{IB 514: Scientific Writing}
\date[]{}

\AtBeginSection[]
{
  \begin{frame}
    \frametitle{Overview}
    \tableofcontents[currentsection]
  \end{frame}
}

\begin{document}

% -------------------------------------------------
\begin{frame}
  \titlepage
\end{frame}

% -------------------------------------------------
\begin{frame}{Session Objectives}
\begin{itemize}
    \item Build sustainable daily writing routines
    \item Design protected writing time
    \item Structure deadlines strategically
    \item Improve drafting and workflow practices
    \item Diagnose and intervene on writer’s block
    \item Develop meta-cognitive awareness
\end{itemize}
\end{frame}

% =================================================
\section{Daily Writing Routines}
% =================================================

\begin{frame}{Daily Writing as a System}
\textbf{Core Principle:} Writing is a repeatable cognitive practice.

\begin{itemize}
    \item Frequency $>$ intensity
    \item Define task before session begins
    \item Protect high-cognition hours
\end{itemize}
\end{frame}

% -------------------------------------------------

\begin{frame}{Some Models of Daily Writing}
\small
\textbf{1. Consistency Model}
\begin{itemize}
    \item 30--60 min daily
    \item Same time, same location
\end{itemize}

\textbf{2. Sprint Model}
\begin{itemize}
    \item 25-min focused intervals
    \item 2--3 per day
\end{itemize}

\textbf{3. Cognitive Load Model}
\begin{itemize}
    \item Morning: argument development, drafting new sections
    \item Midday: data checks, figure refinement
    \item Late day: formatting, references, minor edits
\end{itemize}

\textbf{4. Paragraph Commitment Model}
\begin{itemize}
    \item 1 paragraph/day
    \item Stop after completion
\end{itemize}
\end{frame}

% -------------------------------------------------

\begin{frame}{Concrete Output and Task Decomposition}
\begin{itemize}
    \item 1 paragraph/day $\rightarrow$ 5 per week
    \item 300 words/day $\rightarrow$ 6,000 per month
    \item Draft Results before Introduction
\end{itemize}

\vspace{0.5em}

\textbf{Task Decomposition}
\begin{itemize}
    \item Instead of “Revise Discussion”
    \item[] $\rightarrow$ Cut 300 words
    \item[] $\rightarrow$ Clarify paragraph 3 claim
    \item[] $\rightarrow$ Add 2 citations on stability theory
\end{itemize}
\end{frame}

% =================================================
\section{Time Management \& Deadlines}
% =================================================

\begin{frame}{Protected Writing Time}

Protected time is:
\begin{itemize}
    \item Scheduled
    \item Visible
    \item Non-negotiable
\end{itemize}

\begin{block}{Treat it as a meeting with yourself.}
\end{block}

\textbf{Implementation Strategies}
\begin{itemize}
    \item Block 60--90 min during peak cognitive hours
    \item Close email + messaging platforms
    \item Define task in calendar description
    \item Begin immediately (no warm-up browsing)
    \item Start with 2--3 blocks/week; increase after consistency
\end{itemize}

\end{frame}

% -------------------------------------------------

\begin{frame}{Major vs. Minor Deadlines}

\begin{block}{Nested Planning Structure}
Major deadlines determine direction.\\
Minor deadlines determine daily behavior.
\end{block}

\textbf{Major Deadlines}
\begin{itemize}
    \item Journal submission
    \item Thesis chapter defense
    \item Circulate full draft to co-authors
\end{itemize}

\textbf{Minor Deadlines}
\begin{itemize}
    \item Draft study system section of Methods by April 1
    \item Figure 2 drafted by Tuesday
    \item Circulate Introduction in 10 days
    \item Draft Abstract by Friday
\end{itemize}

\end{frame}

% -------------------------------------------------

\begin{frame}{Examples of Minor and Micro Deadlines}

Examples within a 4-week revision window:

\begin{itemize}
    \item Day 3: Rewrite Results subsection 1
    \item Day 5: Reduce Discussion by 400 words
    \item Week 2: Complete robustness checks
    \item Week 3: Internal logic review of Introduction
    \item Week 4: Format references
\end{itemize}

\vspace{0.5em}

\textbf{Micro-Deadlines}
\begin{itemize}
    \item Clarify claim of paragraph 4 (30 min)
    \item Verify statistical assumption 2
    \item Insert 3 missing citations
    \item Rewrite topic sentences in Results
\end{itemize}

Minor goals should be specific, time-bound, and measurable.

\end{frame}

% =================================================
\section{Workflow Organization}
% =================================================

\begin{frame}{Workflow as a Pipeline}

\begin{block}{
Ideas $\rightarrow$ Notes $\rightarrow$ Analysis $\rightarrow$ Figures $\rightarrow$ Manuscript
}
Goal: reduce switching costs.
\end{block}



\textbf{Explicit Transition Points}
\small
\begin{itemize}
    \item Notes $\rightarrow$ Analysis: Research question formalized
    \item Analysis $\rightarrow$ Figures: Model validated + assumptions checked
    \item Figures $\rightarrow$ Manuscript: Figure narrative drafted in bullet form
    \item Manuscript $\rightarrow$ Circulation: Logical coherence pass completed
\end{itemize}

\textbf{Definition of Done Examples}
\begin{itemize}
    \item Analysis complete = validated + documented
    \item Figure complete = caption explains inference
    \item Section draft complete = each paragraph has explicit claim
\end{itemize}

\end{frame}

% -------------------------------------------------

\begin{frame}{Cognitive Energy Management}

High-energy tasks:
\begin{itemize}
    \item Argument construction
    \item Structural revision
\end{itemize}

Medium-energy tasks:
\begin{itemize}
    \item Sentence rewriting
\end{itemize}

Low-energy tasks:
\begin{itemize}
    \item Formatting
    \item Reference checks
\end{itemize}

Align task to cognitive state.

\end{frame}

% =================================================
\section{Early Drafting Practices}
% =================================================

\begin{frame}{Draft Before You Feel Ready}

\begin{block}{Write imperfect prose intentionally}
    "Write drunk, edit sober."
\end{block}

Separate drafting from editing.

\begin{itemize}
    \item Draft full section quickly
    \item Revise next day
\end{itemize}

\end{frame}

% -------------------------------------------------

\begin{frame}{Reverse Outline in Bullet Points}

\begin{block}{Reverse outlining}
Extracting the logical structure after drafting.
\end{block}

\textbf{Process}
\begin{itemize}
    \item Reduce each paragraph to one bullet:
    \item[] “What claim does this paragraph make?”
    \item Examine logical order.
\end{itemize}

\textbf{Example}

Paragraph claim: \\
{\small \textit{Species diversity increases stability under environmental fluctuation.}}

Reverse outline bullet: \\
{\small \textit{Diversity buffers variance via asynchronous population responses.}}

\vspace{0.5em}
If bullets do not form a coherent argument, restructure.

\end{frame}

% =================================================
\section{Writer's Block}
% =================================================


% -------------------------------------------------

\begin{frame}{Diagnosing Writer's Block: Underlying Causes}

Writer’s block is usually a symptom, not a primary problem.

\vspace{0.5em}

\textbf{1. Cognitive Causes}
\begin{itemize}
    \item Argument is not fully specified
    \item Paragraph lacks a clear claim
    \item Logical steps between ideas are missing
    \item Methods or assumptions not fully understood
\end{itemize}

\textbf{2. Structural Causes}
\begin{itemize}
    \item No outline or unstable outline
    \item Section scope too broad
    \item Unclear transition between sections
\end{itemize}


\end{frame}

\begin{frame}{Diagnosing Writer's Block: Underlying Causes}

Writer’s block is usually a symptom, not a primary problem.

\vspace{0.5em}

\textbf{3. Process Causes}
\begin{itemize}
    \item Editing while drafting
    \item Task too large (“Revise paper”)
    \item No defined next action
\end{itemize}

\textbf{4. Psychological Causes}
\begin{itemize}
    \item Perfectionism
    \item Fear of evaluation (advisor, reviewers)
    \item Avoidance of cognitively demanding work
\end{itemize}

\vspace{0.5em}

\begin{block}{Writer's Block}
A signal that something specific needs clarification.
\end{block}

\end{frame}


% -------------------------------------------------

\begin{frame}{Interventions: Matching Strategy to Cause}

\textbf{If the argument is unclear:}
\begin{itemize}
    \item Write the paragraph’s claim as a single sentence.
    \item Ask: What must be true for this claim to hold?
    \item Sketch logic on paper before drafting prose.
\end{itemize}

\textbf{If the task is too large:}
\begin{itemize}
    \item Reduce to a 15-minute goal.
    \item Replace “Revise Discussion” with
    “Rewrite paragraph 2 for clarity.”
\end{itemize}

\textbf{If perfectionism is interfering:}
\begin{itemize}
    \item Draft in bullet points.
    \item Use timed freewriting (10 minutes, no deletion).
    \item Write a deliberately “bad” version.
\end{itemize}


\end{frame}

\begin{frame}{Interventions: Matching Strategy to Cause}


\textbf{If fatigue is the issue:}
\begin{itemize}
    \item Switch to a lower-cognitive task (references, figure labels).
    \item Stop and resume during peak cognition.
\end{itemize}

\textbf{If stuck mid-paragraph:}
\begin{itemize}
    \item Insert bracketed placeholders:
    [CITE], [EXPLAIN MECHANISM], [CHECK RESULT]
\end{itemize}

\vspace{0.5em}
Momentum restores clarity more reliably than rumination.

\end{frame}


% =================================================
\section{Meta-Cognition}
% =================================================

\begin{frame}{Meta-Cognitive Writing}

Strong writers monitor:
\begin{itemize}
    \item When they write best
    \item How long tasks take
    \item What triggers avoidance
    \item What improves clarity
\end{itemize}

\end{frame}

% -------------------------------------------------

\begin{frame}{Practical Monitoring Systems}

\textbf{Writing Log}
\begin{itemize}
    \item Date
    \item Duration
    \item Words added/removed
    \item Task completed
\end{itemize}

\textbf{Example Log Entry}
\begin{itemize}
    \item 8:00--8:45 AM
    \item Drafted 2 Results paragraphs
    \item Stalled at interpretation step
    \item Insight: need clearer model summary before Discussion
\end{itemize}

\end{frame}

% -------------------------------------------------

\begin{frame}{Closing Reflection}

Write down:
\begin{itemize}
    \item One protected writing block
    \item One minor deadline for this week
    \item One workflow transition to clarify
    \item One drafting experiment
\end{itemize}

\vspace{1em}

\textit{Writing productivity is designed, not discovered.}

\end{frame}

% -------------------------------------------------

\end{document}